\documentclass[10pt,a4paper]{amsart}
\usepackage[margin=19mm]{geometry}
\usepackage{amsmath,amssymb,mathtools}
\usepackage[scr=rsfs,cal=euler]{mathalfa}
\usepackage{xfrac}
\usepackage[unicode,hidelinks,bookmarks=false]{hyperref}
\newcommand{\ie}{i.e.\ }
\newcommand{\cf}{cf.\ }
\newcommand{\resp}{respectively\ }
\newcommand{\Subsection}{\subsection}
\providecommand{\nobreakdash}{\nobreak\hbox{-}}
\setlength{\emergencystretch}{2em}
\multlinegap0pt
\usepackage{bm}
\usepackage{bbm}
\usepackage{dsfont}
\usepackage{xspace}
\usepackage{cancel}
\usepackage{mathtools}
\usepackage{amsthm}
\makeatletter
\@ifundefined{lemma}{\newtheorem{lemma}{Lemma}}{}
\makeatother
\title[On certain combinatorial expressions of TASEP transition pro]{On certain combinatorial expressions of TASEP transition probabilities}
\author{Lorenzo Vito Dal Zovo}
\date{Source: arXiv:2605.29520v1 (CC BY 4.0)}
\begin{document}
\maketitle

\noindent\emph{Source and licence.} On certain combinatorial expressions of TASEP transition probabilities. Version arXiv:2605.29520v1 (CC BY 4.0), distributed under the Creative Commons Attribution 4.0 International licence, as recorded in the arXiv licence metadata for this exact version and shown on the abstract page https://arxiv.org/abs/2605.29520v1. Full rights notice: licenses/arxiv-2605.29520-CC-BY-4.0.txt.\par\smallskip

\noindent\textit{Editorial scope.} Counted unit: the Lemma whose source label is \texttt{None} in the article (source0.tex lines 744--755, source label \texttt{None}), delivered together with the article's own complete original proof of it (source0.tex lines 757--859, one \texttt{proof} environment of 103 lines). No mathematics is written, repaired or re-derived: statement and proof are the article's own text line for line. Required context retained uncounted: none. Every definition, display and label of those lines is retained unchanged. Omitted from this package and not redistributed: 1--743, 756--756, 860--1042. Nothing inside a retained excerpt is omitted or altered: every retained body is delivered exactly as the article's lines are written, so no omission receipt of any kind accompanies this package. Citations: the retained excerpts carry no citation command at all, so no bibliography entry of the article is transcribed and no reference is redistributed. Reference adaptation: none was needed and none was made; every reference inside the delivered unit resolves inside it, so no unresolved reference or citation is produced. Printed slips kept literal: no printed word, symbol, hypothesis or number of the counted statement or of its proof was altered. Layout and class: the article's own class is \texttt{article} with packages including graphicx, amsmath, amsfonts, amssymb, amsthm, hyperref, mathtools, xcolor, tikz, anyfontsize; the wrapper is the lane's pinned \texttt{amsart} editorial wrapper at 10pt on A4, which restates the theorem-like environments the retained text needs and adds the article's own macro definitions that the retained text needs (none), with extra wrapper packages (bm, bbm, dsfont, xspace, cancel, mathtools). The delivered numbering is the unit's own. Assets: no retained span contains a figure environment, a graphics-inclusion command, a PDF-page inclusion, a table or a TikZ picture; no image file is copied into this package, the package declares no asset dependency and no image file is redistributed with it. Licence: version arXiv:2605.29520v1 of this article is distributed under the Creative Commons Attribution 4.0 International licence, as recorded in the arXiv licence metadata for this exact version and shown on the abstract page https://arxiv.org/abs/2605.29520v1. A full rights notice is delivered at licenses/arxiv-2605.29520-CC-BY-4.0.txt. Characters: every retained excerpt is pure ASCII, so the delivered engine, XeLaTeX with its default Latin Modern text font, reports no missing character.
\begin{lemma}
Let D be a diagram in $D_{\xi,\eta}^{N+1}$, where $\xi,\eta\in\{\uparrow,\rightarrow\}^N$, be such that there exists $D_0\in D_{\xi,\eta}^{N+1}$ with $D_0\subset D$. Let $s\in\Sigma^n$ be a word in the equivalence class induced by $(D,\Sigma,\ell_D)$, where $\ell_D$ is the restriction of $\overline{\ell}$ to $D$. Set $x=g(\xi)\in X_N$ and $y=g(\eta)\in X_N$. Then
\[
(s)_{ij}=
\begin{cases}
(-1)^{\,n-|D|+\kappa(z,y)}
& \text{if } j=\operatorname{idx}(x)\ \text{and}\ i=\operatorname{idx}(z)\ \text{for some } z\in \mathcal{N}^-(y),\\[4pt]
0 & \text{otherwise},
\end{cases}
\]
where $\kappa(z,y)$ denotes the number of particles that are moved one site to the left in passing from $y$ to $z$, that is, $\kappa(z,y)=\frac12\|y-z\|_1$).
\end{lemma}

\begin{proof}

We decompose each local generator $h_r\in\Sigma$ as the sum of two matrices, one containing its diagonal component and the other its off diagonal one:
\begin{equation}
    \label{expansion}
    h_r=-h_r^{(1)}+h_r^{(2)},
\end{equation}
where for $1<r<N+1$ we have
\[
h_r^{(1)}=I^{\otimes (r-2)}\otimes(E_{22}\otimes E_{11})\otimes I^{\otimes (N-r)},
\qquad
h_r^{(2)}=I^{\otimes (r-2)}\otimes(E_{12}\otimes E_{21})\otimes I^{\otimes (N-r)},
\]
while at the boundaries
\[
h_{1}^{(1)}=E_{11}\otimes I^{\otimes (N-1)},\qquad
h_{1}^{(2)}=E_{12}\otimes I^{\otimes (N-1)},
\]
and
\[
h_{N+1}^{(1)}=I^{\otimes (N-1)}\otimes E_{22},\qquad
h_{N+1}^{(2)}=I^{\otimes (N-1)}\otimes E_{12}.
\]

Let $\overline{s}$ be a reduced word for $s$. Since each cancellation $h_r^2=-h_r$ removes one letter and contributes one factor $-1$, we may write
\[
\overline{s}=(-1)^{n-m}\prod_{q=1}^{m}\overline{s}_q,
\qquad m=|D|.
\]
Expanding every factor according to \eqref{expansion} gives
\begin{equation}
\label{eq:expansion_lemma}
\overline{s}=(-1)^{n-m}\prod_{q=1}^{m}\bigl(\overline{s}_q^{(2)}-\overline{s}_q^{(1)}\bigr).
\end{equation}
Thus $s$ is a signed sum of $2^m$ Kronecker products, one for each choice of diagonal/off-diagonal contribution at every factor of the reduced word.

We claim that a summand in \eqref{eq:expansion_lemma} is nonzero if and only if the position, relative to the partial order induced by $\overline{s}$, of every factor chosen as a diagonal part corresponds to a corner of the diagram $D$. Indeed, if a diagonal term is chosen at a cell that is not a corner, then one of its neighboring factors in the reduced word leads to a null product, and the resulting term in the expansion vanishes. By contrast, at any corner the diagonal choice is compatible with the surrounding factors and therefore may contribute a non-null term.

Consider first the summand obtained by choosing the off-diagonal part at every factor:
\[
s_0\coloneqq (-1)^{n-m}\prod_{q=1}^{m}\overline{s}_q^{(2)}
      \coloneqq(-1)^{n-m}\bigotimes_{r=1}^{N}c_r.
\]
The factor $c_r$ is determined by the $r$-th step of the two outlines $\xi$ and $\eta$:
\begin{equation}
\label{eq:ci_from_outlines}
c_r=
\begin{cases}
E_{11} & \text{if } \xi_r=\eta_r=\rightarrow,\\
E_{22} & \text{if } \xi_r=\eta_r=\uparrow,\\
E_{12} & \text{if } \xi_r=\rightarrow,\ \eta_r=\uparrow,\\
E_{21} & \text{if } \xi_r=\uparrow,\ \eta_r=\rightarrow.
\end{cases}
\end{equation}
On the other hand, the basis matrix $E_{\operatorname{idx}(x),\operatorname{idx}(y)}$ admits the factorization
\begin{equation}
\label{eq:ci_from_configs}
E_{\operatorname{idx}(y),\operatorname{idx}(x)}=\bigotimes_{r=1}^{N}d_r,
\qquad
d_r=
\begin{cases}
E_{11} & \text{if } x_r=y_r=0,\\
E_{22} & \text{if } x_r=y_r=1,\\
E_{12} & \text{if } x_r=0,\ y_r=1,\\
E_{21} & \text{if } x_r=1,\ y_r=0.
\end{cases}
\end{equation}
Since the map $g$ identifies $\uparrow$ with $1$ and $\rightarrow$ with $0$, the descriptions in \eqref{eq:ci_from_outlines} and \eqref{eq:ci_from_configs} coincide. Therefore
\[
s_0=(-1)^{n-m}E_{\operatorname{idx}(y),\ \operatorname{idx}(x)}.
\]

Now choose a nonempty set $C$ of corners of $D$, and in \eqref{eq:expansion_lemma} select the diagonal part exactly at the factors whose relative order is given by the cells in $C$ and call $s_1$ the corresponding term in the expansion \eqref{eq:expansion_lemma}. Because factors whose relative position is given by corners commute with each other, the corresponding summand is again a nonzero Kronecker 
product with a single nonzero entry. For the same reason, the Kronecker factorization of $s_1$ corresponds to that of $s_0$ apart from the factors corresponding to the cells in $C$. Let $i,\,i+1$ be the indices associated to a corner in $C$ and $b_i,\, b_{i+1}$ the respective Kronecker factors of $s_1$, then:
 \begin{equation*}
    b_{i}=\left\{
        \begin{array}{lll}
         E_{12}  &\text{if }  &\eta_{i}=\xi_{i}=\uparrow, \\ 
         E_{11}  &\text{if }  &\eta_{i}=\uparrow, \ \xi_{i}=\rightarrow,
        \end{array}\right.
\end{equation*}
and
\begin{equation*} 
    b_{i+1}=\left\{
        \begin{array}{lll}
         E_{21}  &\text{if }  &\eta_{i+1}=\xi_{i+1}=\rightarrow, \\ 
         E_{22}  &\text{if }  &\eta_{i+1}=\rightarrow, \ \xi_{i+1}=\uparrow,
        \end{array}\right.
    \end{equation*}

Thus, relative to $s_0$, each chosen corner replaces one allowed right jump by a "stay" at the corresponding particle position; on the level of configurations this means that the arrival configuration $y$ is modified by moving that particle one site to the left. Hence the resulting term is
\[
(-1)^{n-m+|C|}E_{\operatorname{idx}(z_C),\operatorname{idx}(x)},
\]
where $z_C\in\mathcal{N}^-(y)$ is obtained from $y$ by moving to the left exactly the particles corresponding to the corners in $C$. Observing that $\kappa(z_C,y)=|C|$, then every nonzero summand of \eqref{eq:expansion_lemma} is of the form
\[
(-1)^{n-|D|+\kappa(z,y)}E_{\operatorname{idx}(z),\operatorname{idx}(x)}
\qquad\text{for some }z\in\mathcal{N}^-(y).
\]

Collecting all summands proves that the only nonzero entries of $s$ occur in column $\operatorname{idx}(x)$ and in the rows indexed by configurations $z\in\mathcal{N}^-(y)$, with coefficient $(-1)^{n-|D|+\kappa(z,y)}$ one recovers the thesis.

\end{proof}

\end{document}
